\section{Conclusion}\label{sec:conclusion}

Family transfers are a substantial part of recorded residential turnover:
nearly 20 million deeds from 2007 through 2023, or one for every 3.6 market
sales. Their incidence differs across states and housing vintages, and
many transferred properties remain unsold for years. These facts establish
the scale and persistence of a channel that price-based transaction data
often omit. They do not establish that family ownership is inefficient.

California's Proposition~19 shows that this channel responds to tax
incentives. The anticipation surge and subsequent decline in family deeds
are consistent with a substantial role for inherited assessments in
transfer decisions. The evidence on later market sales is less decisive:
there is no statistically clear increase in estate/trust-seller sales
through 2023, and the conditional sale-hazard estimates are imprecise.
Deferral is one explanation, but altered recording and other responses
cannot be excluded by those null results.

The framework separates efficient family retention from retention induced
by a tax advantage. Under explicit assumptions about allocation response,
tax wedges, succession delays, and persistence, the calibration translates
the deed response into potential property-years and dollars of allocation
gain. It also shows why the same deed response is compatible with zero
gains: legal ownership can change without housing use changing, and an
inherited rental can continue to serve a tenant efficiently. Tax revenue
and housing surplus are separate accounts.

Longer follow-up and transaction-time information on exclusions, occupancy,
rentals, and assessments could identify which missing deeds correspond to
different allocations and narrow the calibration. The policy question is
not whether families should retain homes. It is whether property-tax inheritance
rules reward retention when another allocation would generate greater net
surplus, and how often that margin is relevant in practice.
